\chapter{Introduction}

\section{Context}

Kolkata's road network carries millions of trips daily across a mix of private
vehicles, metros, busses, and pedestrians. Disruptions -- waterlogging, VIP
convoys, accidents, road closures -- are frequent and hard to predict from
sensor data alone because a large share of disruption information first appears
in news articles and official advisories, not on roads.

Existing navigation tools (Google Maps, Apple Maps) react to congestion after
it shows up in GPS probe data, often minutes to hours after the event begins.
This project asks whether reading that unstructured text \textit{before travel}
can produce better route recommendations.

\section{Problem Statement}

\begin{table}[H]
\centering

\begin{tabular}{lll}
\toprule
\textbf{Dimension} & \textbf{Current systems} & \textbf{This project} \\
\midrule
Data types used     & Structured only          & Structured + unstructured \\
Timing              & After disruption         & Before disruption \\
Objective           & Minimise travel time     & Minimise risk-adjusted time \\
Uncertainty         & Deterministic            & Probabilistic \\
Explainability      & None                     & HGNN attention weights \\
Multimodal          & Limited                  & 8 modes inc.\ metro and bus \\
\bottomrule
\end{tabular}
\caption{Reactive vs.\ anticipatory trip planning}
\end{table}

\section{Research Hypotheses}

\begin{table}[H]
\centering

\begin{tabular}{cp{11.5cm}}
\toprule
\textbf{ID} & \textbf{Hypothesis} \\
\midrule
H1 & LLM-derived contextual signals improve disruption detection accuracy
     over structured-data-only baselines. \\
H2 & Probabilistic, risk-aware routing improves travel-time reliability
     under uncertain conditions. \\
H3 & Exposing HGNN attention weights as route explanations increases
     commuter trust in system recommendations. \\
\bottomrule
\end{tabular}
\caption{Guiding hypotheses}
\end{table}

\section{Scope of Internship}

\begin{table}[H]
\centering
\caption{Work carried out during the internship}
\begin{tabular}{clc}
\toprule
\textbf{\#} & \textbf{Contribution} & \textbf{Status} \\
\midrule
1  & LLM disruption extraction pipeline (Layer 1)               & Done \\
2  & Multi-factor confidence scoring + impact duration           & Done \\
3  & Temporal HGNN -- training loop + inference                 & Done \\
4  & Bayesian probabilistic fusion (Layer 2)                     & Done \\
5  & Weather Severity Index per route                            & Done \\
6  & Kolkata Metro integration (5 lines, timetable, next-train)  & Done \\
7  & Bus transit graph + BFS routing engine                      & Done \\
8  & Multimodal routing orchestration                            & Done \\
9  & FastAPI backend (11 endpoints)                              & Done \\
10 & Leaflet.js frontend + map overlays                          & Done \\
11 & CVaR risk-aware routing (Layer 3)                           & In progress \\
\bottomrule
\end{tabular}
\end{table}

\section{Report Structure}

\begin{table}[H]
\centering
\begin{tabular}{cl}
\toprule
\textbf{Chapter} & \textbf{Content} \\
\midrule
2 & System architecture, data sources, API design \\
3 & Implementation -- Layers 1, 2, routing, database \\
4 & Results, UI screenshots, performance metrics \\
5 & Conclusion and future work \\
\bottomrule
\end{tabular}
\end{table}
